\subsection*{Что даёт опыт при выборе эталона}
На изотерме 1273,15~К CALPHAD отличается от измерений Кайнумы на
\(-0{,}15\) п.п. для \(\alpha\) и \(+0{,}19\) п.п. для \(\alpha_2\).
Оба расхождения укладываются в приведённую авторами точность.
При 1323,15~К разности составляют \(-0{,}56\) и \(+0{,}46\) п.п.
Близость экспериментальных и расчётных границ зависит от температуры,
поэтому усреднение соседних точек вокруг 1300~К здесь скрыло бы часть
измеренного изменения.

На прежнем проектном рисунке точки Блэкберна оцифрованы с
рис.~13 Николса и соавт. Вблизи 1000~\(^\circ\)C они дают примерно
19 и 24~ат.\% Al. Для численного сравнения выше использована исходная
таблица Кайнумы. Николс и соавторы также исследовали переход на MTP
по скачку состава в двух сериях при заданных химических полях,
с начальными структурами \(\alpha\) и D0$_{19}$. Для MTP выше
1100~К петля гистерезиса стала слишком мала для определения границ
этим способом. Порог наблюдаемости петли зависит от метода и времени
наблюдения, поэтому при 1273,15~К мы отдельно проверили относительные
равновесные веса.
\href{https://doi.org/10.3390/met16020140}{Nichols et al. (2026), разд.~2.4, табл.~2, рис.~13.}
